\section{Related Work}

Our work connects three research threads: sub-quadratic sequence models, efficient adaptation methods, and architecture conversion. Each thread has made significant progress independently, but none addresses the adaptation preservation problem during conversion.

\subsection{Sub-Quadratic Sequence Models}

The quadratic complexity of transformer attention has motivated extensive work on efficient alternatives. Longformer \cite{Beltagy2020Longformer} and BigBird \cite{Zaheer2020BigBird} introduce sparse attention patterns that reduce complexity to $O(n)$. Linear attention methods \cite{Katharopoulos2020Transformers} reformulate attention as a kernel feature map, achieving $O(n)$ complexity but sacrificing expressiveness.

More recently, state space models have emerged as a principled sub-quadratic alternative. Mamba \cite{Gu2023Mamba} introduces selective state spaces with input-dependent gating, achieving transformer-level performance at $O(n)$ complexity. RWKV \cite{Peng2023RWKV} combines linear attention with RNN-style recurrence. These architectures excel at efficiency but provide no mechanism for task-specific adaptation during deployment. Our work extends Mamba with task conditioning that preserves adaptation capability.

\subsection{Efficient Adaptation Methods}

Adapting large models to downstream tasks efficiently has driven development of parameter-efficient methods. LoRA \cite{Hu2021LoRA} learns low-rank updates to pretrained weights, reducing trainable parameters to $\sim$0.1\% while maintaining performance. Adapters \cite{Houlsby2019Adapters} insert small bottleneck modules between transformer layers. Prompt tuning \cite{Lester2021Prompt} optimizes soft prompts prepended to inputs.

These methods assume the base model preserves adaptation capability---an assumption that breaks when converting architectures. LoRA applied post-conversion must work with degraded internal representations. TC-SSM differs fundamentally: we integrate task conditioning \textit{during} conversion, preserving the adaptation manifold rather than compensating for its loss.

\subsection{Model Conversion and Distillation}

Knowledge distillation \cite{Hinton2015Distilling} transfers knowledge from teacher to student models by matching output distributions. For architecture conversion, distillation typically optimizes KL divergence on logits and MSE on hidden states. StreamingLLM \cite{Xiao2023StreamingLLM} and H2O \cite{Zhang2023H2O} focus on KV cache compression for efficient inference, maintaining output quality without changing the architecture.

The conversion literature optimizes for output fidelity, not adaptation preservation. A converted model may match teacher accuracy on fixed tasks while losing the ability to quickly adapt to new ones. TC-SSM introduces an adaptation regularizer during conversion training that explicitly preserves task-conditioned dynamics, addressing a gap that pure output matching cannot fill.

\subsection{Positioning TC-SSM}

Prior work addresses efficiency (Mamba, RWKV), adaptation (LoRA, adapters), or output preservation (distillation) in isolation. TC-SSM is the first method to co-design conversion and adaptation, integrating task conditioning into the conversion process itself. This enables sub-quadratic models that retain transformer-level few-shot capability without post-hoc compensation.
